\section{Introduction}
Misinformation spreads through short, image–rich posts where cues in text and vision interact in subtle ways. Multimodal detectors improve over text–only systems but remain vulnerable to topic shift and annotation noise; the literature explores event–invariant representations \citep{Wang2018EANN}, similarity–aware fusion \citep{Zhou2020SAFE}, graph/social context \citep{Lu2020GCAN,Nguyen2020FANG}, and large–scale benchmarks such as Fakeddit \citep{Nakamura2020Fakeddit}. In parallel, the vision–language community has shifted toward pretrained encoders (e.g., ViLBERT, CLIP) that better capture cross–modal semantics \citep{ViLBERT2019,Radford2021CLIP}, yet deploying them robustly on domain–specific misinformation remains challenging.

Knowledge distillation (KD) transfers supervision from accurate or expensive teachers to a compact student \citep{Hinton2015Distillation,Gou2021KDsurvey}, including mutual/self–distillation variants \citep{Zhang2018DML,Furlanello2018BAN}. However, in noisy, shifting domains like misinformation detection, teachers can be wrong with high confidence; naively imitating them propagates errors and degrades calibration. Robust learning under noise selects reliable instances or mentors during training \citep{Han2018Coteaching,MentorNet2018}, but comparable \emph{selective} mechanisms for distillation from \emph{multimodal teachers} are under–explored.

Uncertainty and calibration are crucial for safety and downstream moderation. Modern neural networks are often miscalibrated \citep{Guo2017Calibration}; evidential learning and prior networks explicitly model predictive uncertainty \citep{Sensoy2018EDL,Malinin2018PN}. For KD, temperature control and cost–sensitive objectives can mitigate over–confident transfer, especially on rare positives that matter most to users.

In this work we focus on a concrete and practically important setting: supervised multimodal rumor detection with \emph{multi-branch teachers and explicit event identifiers}. We assume a fixed multimodal teacher with text, vision, and fusion branches (as in EANN/SAFE-style encoders) and event IDs for each post (rumor/story IDs in Weibo; public-account stories in WeFEND). Our goal is not a universal KD framework, but to study how selective, reliability-aware transfer from such teachers can reduce error propagation under event-structured noise.

We propose \textbf{Dynamic Modal Priority Distillation (DMPD)}, a \emph{reliability-aware selective distillation} scheme for \taskname{}. DMPD routes supervision from a multimodal teacher with text, vision, and fusion branches to the student only when (i) a \emph{teacher–match gate} confirms cross-modal agreement and confidence, and (ii) an \emph{event–level reliability} score indicates consistent teacher correctness. A selective distillation objective adds safety constraints—temperature–scaled KL, lightweight feature alignment, positive-class costs, and evidential regularization—to avoid over-confident transfer while preserving recall on the positive (fake) class. Compared to classical KD that copies teacher outputs indiscriminately, DMPD performs agreement- and reliability-aware selection plus calibration to limit error propagation under structured (event-level) noise. We use “DMPD” henceforth.

Our contributions (all serving the single goal of reliability-aware selective distillation) are:
\begin{itemize}
  \item \textbf{Selective gating under event-structured noise:} a lightweight teacher–match gate and an event-reliability EMA that jointly admit only cross-view-consistent, historically reliable supervision, suppressing bursty or biased events.
  \item \textbf{Safety constraints for selective KD:} a calibration-oriented objective (temperature-annealed KL, evidential head) and positive-class costs that keep gains from selective distillation without over-confident transfer.
  \item \textbf{Targeted empirical study:} experiments on two multimodal rumor datasets with explicit event IDs (Weibo, WeChat/WeFEND), multi-seed statistics, and ablations on gating/reliability; plus an English synthetic OOD benchmark (MiRAGeNews) and a strict time-split Fakeddit “hard case,” showing where selective distillation helps and where weak teachers limit gains.
\end{itemize}

Comprehensive experiments, ablations, and visualizations are detailed in Sections~\ref{sec:experiments}--\ref{sec:analysis}. We focus on settings where a reasonably strong multimodal teacher is available and aim to make distillation safer and better calibrated in the presence of event-level noise, rather than rescuing arbitrarily weak teachers. Code, configs, and scripts will be released for full reproducibility once anonymity is lifted.
